\scanpage{20}
puis
\[
\rho^3(\bar a)=\bar a
\]
Donc
\[
\rho^3=\mathrm{id}_{\bar A}
\]
En particulier, $\rho$ est une \uline{permutation}. Toutes ses orbites sont de cardinal 3. Donc il y a cinq telles orbites.

Une orbite de $\rho$ s'appelle \uline{trièdre} \uline{distingué} \inserted{(de directions d'arêtes)} associé à $\rho$. Deux directions d'arêtes \inserted{$\alpha,\beta$} sont dites \uline{orthogonales} si elles sont \inserted{$\ne$ et} dans la $=$ orbite, \inserted{i.e.  \deleted{\unsure{arête}} $=$ trièdre distingué}, ce qui signifie aussi
\[
\beta=\rho\alpha\quad\text{ou}\quad\alpha=\rho\beta
\]
Pour tte \deleted{orthog} direction d'arête, il y a donc exactement deux directions d'arêtes orthogonales. [On notera que la relation d'orthogonalité est sym. par définition, mais la relation $\beta=\rho\alpha$ \deleted{\unsure{$\alpha$}} \inserted{entre $\alpha$ et $\beta$} ne l'est pas, évidemment.\unsure{]}

L'ens des trièdres distingués est un ens quotient de cardinal $5$ de $\bar A$, noté \deleted{\unsure{$T$} \unsure{$I$}} $T$. \inserted{\deleted{\unsure{$A$}} On a application canonique $\bar A\longrightarrow T$} On peut le considérer comme un ens quotient de $A$, formé de $5$ parties de cardinal $6$, chacune d'elles stable par $\underline a$.
% slanted
\margin{terminologie : arêtes orthogonales}

\deleted{\unsure{In}}Interprétation concrète en termes de coloriages,
% archive sheet number: 19
